\section{An actual partial return to the smooth Meyer quotient}
\label{sec:partial-return}

The obstruction in MA-HM-003 concerns a return that recovers
every smooth spectral block.  There is a different return to
the closed smooth subrepresentation supported on the zeros that
the Hilbert completion retains.  At generic widths, compact
Gevrey smoothing makes this return nuclear.  The construction
below uses actual Mellin restrictions and entire interpolation
maps; comparison of scalar zero sums is not used to establish
its operator identities.

Fix the primitive unitary Hecke function and notation of
MT-M18 and MA-WP.  To distinguish the two quotients, write
\[
\begin{aligned}
 Q_\infty&=\mathcal E/\mathcal I,\qquad q_\infty:\mathcal E\to Q_\infty,
       &\mathcal I&=\overline{\psi*\mathcal E}^{\,\mathcal E},\\
 H&=L^2(\mathbb R,w_{u,v,s}(x)\,dx),&
 w_{u,v,s}(x)&=\max(e^{2ux},e^{2vx})(1+x^2)^{s/2},\\
 J_H&=\overline{T_a\mathcal I}^{\,H},&
 Q_H&=H/J_H,\qquad \pi:H\to Q_H,\\
 i:Q_\infty&\longrightarrow Q_H,&
 i[q_\infty f]&=\pi(T_af),\qquad T_af(x)=e^{ax}f(x).
\end{aligned}
\tag{MA-PR.1}
\]
Here $u<v$, $s\ge0$, and all constants concern these fixed
parameters and this fixed Hecke function.  The map $i$ is
continuous and has dense image, but it need not be injective.
Let $\rho=\beta+i\gamma$ be the distinct zeros of the entire
pole-removed completion $\Xi$, with multiplicities $m_\rho$.

Assume the following quantitative avoidance of both edges:
\[
 \min\{|\beta-(a+u)|,\ |\beta-(a+v)|\}
       \ge c_0(2+|\gamma|)^{-2}
       \qquad(\Xi(\rho)=0)
\tag{MA-PR.2}
\]
for one $c_0>0$.  This excludes boundary zeros.
For the fixed function, it holds for almost every pair
$(u,v)$ with $u<v$.  For centered intervals
$a=1/2$, $u=-\alpha$, $v=\alpha$, it holds for almost every
$\alpha>0$.  Indeed the local zero count makes
$\sum_\rho(2+|\gamma|)^{-2}<\infty$; apply the
Borel--Cantelli argument of MA-RS-001 to both edge
parameters and remove the countable set of exact equalities.
For centered intervals the two families of forbidden intervals
are centered at $\beta-1/2$ and $1/2-\beta$.
No zero-location assertion is required.

Set
\[
\begin{aligned}
 Z_{\rm in}&=\{\rho:a+u<\operatorname{Re}\rho<a+v\},\\
 Z_{\rm out}&=\{\rho:\operatorname{Re}\rho<a+u
                         \text{ or }\operatorname{Re}\rho>a+v\},\\
 Q_{\rm in}
   &=\{[f]\in Q_\infty:
       (\mathcal Mf)^{(j)}(\rho)=0
       \text{ for }\rho\in Z_{\rm out},\ 0\le j<m_\rho\}.
\end{aligned}
\tag{MA-PR.3}
\]
This is a closed, translation-invariant subspace of the actual
smooth quotient.  Define $Q_{\rm out}$ by the same formula with
$Z_{\rm in}$ and $Z_{\rm out}$ interchanged.  These definitions
do not assert a continuous unsmoothed projection onto either
subspace.

\paragraph{MA-PR-001 (inner contour restrictions and partial entire lifts).}
There are finite-dimensional Banach algebras $J_n^{\rm in}$,
continuous maps
\[
 r_n:Q_H\longrightarrow J_n^{\rm in},
 \qquad \mathcal T_n^{\rm in}:J_n^{\rm in}\longrightarrow\mathcal E,
\tag{MA-PR.4}
\]
and constants, with $\nu_n=2+|n|$, such that
\[
\begin{gathered}
 d_n^{\rm in}:=\dim J_n^{\rm in}=O(\log\nu_n),\qquad
             \|r_n\|\le C\nu_n,\\
 p_{N,j}(\mathcal T_n^{\rm in}b)
        \le C_{N,j}\exp(C_{N,j}\log^2\nu_n)\|b\|_n,\\
 r_m\,i\,q_\infty\mathcal T_n^{\rm in}
        =\delta_{mn}\operatorname{id}_{J_n^{\rm in}}.
\end{gathered}
\tag{MA-PR.5}
\]
The smooth class of every lift has exactly its prescribed
full jets in the interior zeros of block $n$, and zero full
jets at all other zeros of $\Xi$.
The restriction $i_{\rm in}:Q_{\rm in}\to Q_H$ is continuous,
injective and has dense image.

\begin{proof}
Use the horizontal heights $h_n$ of MT-M18, with
$n+1/4<h_n<n+3/4$.  Choose a sufficiently small fixed
$\delta_0>0$ and put
\[
 \delta_n=\delta_0\nu_n^{-2},\qquad
 \Omega_n^{\rm in}
   =\{a+u+\delta_n<\operatorname{Re}z<a+v-\delta_n,\
                          h_n<\operatorname{Im}z<h_{n+1}\}.
\tag{MA-PR.6}
\]
Choose $\delta_0<(v-u)/4$ and $\delta_0\le1$.
In each fixed-height window about block $n$, (MA-PR.2)
bounds the distance of every zero from either original
vertical edge below by $c\nu_n^{-2}$, with $c$ independent
of $n$.  Make $\delta_0$ smaller than one quarter of this
constant.  The two inner edges then stay at distance at
least $c'\nu_n^{-2}$ from every zero in that window.
The rectangle contains exactly the zeros of $Z_{\rm in}$
whose ordinates are between $h_n$ and $h_{n+1}$.
Its boundary contains no zero.  Its width stays bounded
below by $(v-u)/2$, its height is between $1/2$ and $3/2$,
and its perimeter is uniformly bounded.

Let $J_n^{\rm in}$ be the algebra $A(\Omega_n^{\rm in})$
modulo its ideal of full jets at those interior zeros,
with the quotient of the boundary supremum norm.
This is the same finite holomorphic quotient norm as
(M18.1), for the new rectangle.  Polynomial Hermite
interpolation proves the finite dimension, and the local
zero count gives $d_n^{\rm in}=O(\log\nu_n)$.
Its norm is submultiplicative.  An empty block is the zero
algebra.

For $h\in H$, the function
$w\mapsto\mathcal Mh(w-a)$ is holomorphic on
$a+u<\operatorname{Re}w<a+v$ and continuous on the
closed inner rectangle.  Cauchy--Schwarz gives, for
$u+\delta_n\le\operatorname{Re}z\le v-\delta_n$,
\[
\begin{aligned}
 |\mathcal Mh(z)|
 &\le\|h\|_H
  \left(\int_{\mathbb R}
       \frac{e^{2\operatorname{Re}z\,x}}
            {\max(e^{2ux},e^{2vx})(1+x^2)^{s/2}}\,dx\right)^{1/2}\\
 &\le C\delta_n^{-1/2}\|h\|_H
 \le C'\nu_n\|h\|_H .
\end{aligned}
\tag{MA-PR.7}
\]
The elementary integral bound uses $s\ge0$ and
$e^{-2\delta_n|x|}$ on the two tails.  Restricting this
representative to the quotient algebra defines a bounded
map $R_n^H:H\to J_n^{\rm in}$ of norm at most $C'\nu_n$.
It vanishes on $T_a\mathcal I$, hence on $J_H$, and
therefore descends to $r_n$ in (MA-PR.4).
For $h=T_af$ it gives the ordinary block jets of
$\mathcal Mf$.  No individual-jet condition number enters
this restriction estimate.

We next establish the exact lower bound needed for the
entire lift.  On the horizontal edges, (M18.8) gives its
polynomial minimum modulus on the fixed real interval
containing all the new rectangles.  On a vertical inner
edge use the finite-disk Blaschke argument of MA-RS-001,
or the proof of (M18.8), centered at $2+it$.
Choose its fixed radius large enough that this entire
real interval lies in the inner quarter disk.
There are $O(\log\nu_n)$ zeros in that disk.
Condition (MA-PR.2) and the choice of $\delta_0$ bound
each normalized Blaschke factor below by
$c\nu_n^{-2}$.  The zero-free factor is bounded below
by a fixed inverse power of $\nu_n$, by Harnack's
inequality and the Euler lower bound at $2+it$.
The upper and lower Stirling bounds for $\Xi/L$ then give
\[
 \inf_{w\in\partial\Omega_n^{\rm in}}|\Xi(w)|
       \ge c\exp(-\kappa|n|-C\log^2\nu_n),
       \qquad \kappa=\pi[K:\mathbb Q]/4.
\tag{MA-PR.8}
\]
Only the fixed-height windows used in the chosen disk
matter in selecting $\delta_0$ above.  At the finitely
many low heights, all the contours avoid zeros and
compactness supplies the bound after decreasing $c$.
Finite gamma exceptions and possible finite-$L$ poles
are absent from the large-height disk argument;
$\Xi$ itself is entire on every contour.

Fix $\tau>0$.  For a representative
$F\in A(\Omega_n^{\rm in})$, define
\[
 H_n^{\rm in}F(z)
   =\frac1{2\pi i}\int_{\partial\Omega_n^{\rm in}}
      F(w)\frac{\Xi(z)-\Xi(w)}{\Xi(w)(z-w)}
                         e^{\tau(z-w)^2}\,dw .
\tag{MA-PR.9}
\]
The divided difference is entire in $z$.
The Cauchy argument of (M18.10)--(M18.11) gives
$H_n^{\rm in}F=F+\Xi A$ inside the rectangle,
where $A$ is holomorphic there, and divisibility by
$\Xi$ in a neighborhood of every exterior zero.
If the interior jets of $F$ vanish, $F/\Xi$ is
holomorphic on the rectangle.  The two exterior
Cauchy integrals then vanish; entireness implies
$H_n^{\rm in}F=0$ everywhere.
Thus (MA-PR.9) depends only on the class $b\in J_n^{\rm in}$
and has precisely the stated full jet values.

For each fixed real strip $|\operatorname{Re}z|\le A$,
the proof of (M18.12), with (MA-PR.8) in place of
(M18.9), gives
\[
 |H_n^{\rm in}b(\sigma+it)|
    \le C_A\exp(C_A\log^2\nu_n)\|b\|_n
                                   e^{-c_\tau(t-n)^2}.
\tag{MA-PR.10}
\]
To spell out its estimates, when $|z-w|\ge1$, the
ratio of the upper bound for $\Xi(z)$ to (MA-PR.8)
is bounded by $\exp(C_A\log^2\nu_n)$ times a fixed
polynomial in $2+|t|$ and $e^{\kappa|t-n|}$.
The factor $e^{\tau(z-w)^2}$ absorbs that polynomial
and linear exponential into a smaller Gaussian.
The second term of the divided difference has the
same bound.  When $|z-w|<1$, integrate $\Xi'$ along
the segment from $w$ to $z$ and use its fixed-strip
Stirling upper bound; division by (MA-PR.8) again
costs only $\exp(C_A\log^2\nu_n)$.
The perimeter is bounded, and taking the infimum over
representatives gives the displayed quotient-norm estimate.
Inverse Mellin transformation, shifting to the two
lines $\operatorname{Re}z=\pm(N+1)$ as in MT-M18,
now defines $\mathcal T_n^{\rm in}b\in\mathcal E$.
The integrals of $(1+|t|)^j$ against that Gaussian
prove the second estimate in (MA-PR.5).
The jet identities prove its last equality.

The restriction of $i$ to $Q_{\rm in}$ is injective:
MA-WP-001 says that its kernel has all interior jets
zero, whereas membership in $Q_{\rm in}$ says all
exterior jets are zero.  MT-M18's full smooth kernel
then makes the smooth class zero.  Its image is dense.
Indeed MA-WP's bounded projections $P_n$ have densely
spanning finite-dimensional ranges, each with exactly
the interior local factors because there are no
boundary zeros.  The lifts in (MA-PR.9) realize every
one of those finite factors in $i(Q_{\rm in})$.
This proves the density assertion.
\end{proof}

\paragraph{MA-PR-002 (nuclear partial smoothing with actual factorization).}
Let $g$ have compact support and finite Gevrey order $q>1$,
and put $k_g(x)=e^{-ax}g(x)$.
There is a canonical continuous map
$S_g:Q_H\to Q_{\rm in}\subset Q_\infty$, nuclear of every
positive order, such that
\[
             i_{\rm in}S_g=C_g^{Q_H},\qquad
             S_g i_{\rm in}=C_{k_g}^{Q_{\rm in}}.
\tag{MA-PR.11}
\]
It has an actual continuous lift to $\mathcal E$.
The nuclear bounds are uniform for test families with
fixed compact support, fixed Gevrey order and fixed
Gevrey derivative constants.

On the full smooth quotient define
\[
       B_g=S_g i,\qquad E_g=C_{k_g}^{Q_\infty}-B_g .
\tag{MA-PR.12}
\]
Both are nuclear of every positive order.
$B_g$ has exactly the smoothed interior full jets and
zero exterior full jets; $E_g$ has exactly the smoothed
exterior full jets and zero interior full jets.
Their ranges are contained in $Q_{\rm in}$ and
$Q_{\rm out}$, respectively.  In particular
\[
\begin{gathered}
 B_gB_h=B_{g*h},\qquad E_gE_h=E_{g*h},\\
 B_gE_h=E_gB_h=0
\end{gathered}
\tag{MA-PR.13}
\]
for compact finite-order Gevrey $g,h$.
No unsmoothed projection is asserted.

\begin{proof}
In $J_n^{\rm in}$ let $M_{g,n}$ multiply by the class
of $w\mapsto\mathcal Mg(w-a)$.
The elementary Gevrey Fourier estimate (MA-RS.18)--(MA-RS.19),
uniform on the fixed real interval of the contours, gives
\[
 \|M_{g,n}\|\le
   \sup_{\partial\Omega_n^{\rm in}}|\mathcal Mg(w-a)|
       \le b_n:=C e^{-c\nu_n^{1/q}} .
\tag{MA-PR.14}
\]
We use a positive upper bound $b_n$, including when the
actual multiplier is zero.  The constants $C,c>0$ can
be chosen uniformly for the stated bounded Gevrey families.
Define
\[
 \widetilde S_g
    =\sum_n\mathcal T_n^{\rm in}M_{g,n}r_n:
                        Q_H\longrightarrow\mathcal E,
       \qquad S_g=q_\infty\widetilde S_g .
\tag{MA-PR.15}
\]
Equations (MA-PR.5) and (MA-PR.14) show that the series
converges absolutely in every $\mathcal E$ seminorm,
uniformly on the Hilbert unit ball:
\[
 p_{N,j}(\mathcal T_n^{\rm in}M_{g,n}r_n v)
 \le C_{N,j}\nu_n
       e^{C_{N,j}\log^2\nu_n-c\nu_n^{1/q}}\|v\|_{Q_H}.
\tag{MA-PR.16}
\]
This proves continuity and gives the claimed actual lift.
Its exterior full jets vanish termwise and by continuity,
so $S_g$ has range in the closed subspace $Q_{\rm in}$.

Here is a common nuclear series, to control the topology
as well as the scalar trace.  Choose an Auerbach basis
$e_{n,j}$ of $J_n^{\rm in}$ with dual basis $\ell_{n,j}$;
both vectors and functionals have norm one.  Take
$R_n=C_1\nu_n$ above $\|r_n\|$.  Then
\[
\begin{aligned}
 u_{n,j}&=b_n^{1/2}\mathcal T_n^{\rm in}e_{n,j},\\
 \alpha_{n,j}(v)&=
       \frac{\ell_{n,j}(M_{g,n}r_nv)}{b_nR_n},\\
 \lambda_{n,j}&=R_n b_n^{1/2},\qquad
 \widetilde S_gv=\sum_{n,j}
                       \lambda_{n,j}u_{n,j}\alpha_{n,j}(v).
\end{aligned}
\tag{MA-PR.17}
\]
The $\alpha_{n,j}$ have norm at most one on $Q_H$.
The family $(u_{n,j})$ is bounded in every seminorm of
$\mathcal E$, because its negative Gevrey exponential
dominates each $\exp(C_{N,j}\log^2\nu_n)$.
The coefficients, repeated $d_n^{\rm in}=O(\log\nu_n)$
times, belong to $\ell^p$ for every $p>0$.
This proves nuclearity of every positive order, with
one common bounded target disk for a bounded Gevrey
family.  The images of this disk under $q_\infty$ lie
in $Q_{\rm in}$ and are bounded for its subspace
topology.  Thus the same assertion holds with that
space as target.

For every $m$, the jet identities of the lifts give
\[
       r_m i S_g=M_{g,m}r_m.
\tag{MA-PR.18}
\]
Convolution on the actual quotient also satisfies
$r_m C_g^{Q_H}=M_{g,m}r_m$.
This follows first on the dense smooth image, from
the Mellin convolution identity, then on $Q_H$ by
boundedness of both maps.
By MA-WP-001, the intersection of the kernels of
all the $r_m$ is zero: together they are exactly all
the retained interior jets.
Consequently (MA-PR.18) proves
$i S_g=C_g^{Q_H}$ as an equality of actual bounded
operators on all of $Q_H$.
No comparison of traces has occurred.

For a class in $Q_{\rm in}$, $S_g i$ has the same
interior full jets as $C_{k_g}$, and both have zero
exterior full jets.  The full smooth jet kernel of
MT-M18 proves the second identity in (MA-PR.11).
For an arbitrary class in $Q_\infty$, the same argument
proves the asserted jet description of $B_g$.
It also shows that $S_g$ is independent, as a map to
the smooth quotient, of the allowed choices of
contours, representatives and lifts.
Alternatively its first identity uniquely determines
it among maps into $Q_{\rm in}$, because $i_{\rm in}$
is injective.

The series (MA-PR.17) composed with the continuous $i$
proves nuclearity of every positive order of $B_g$.
The full smooth convolution $C_{k_g}^{Q_\infty}$
has that property by the arithmetic block expansion
of MT-M18--M19.  Their difference $E_g$ therefore
has it too.  Its jet description follows from
subtraction, and gives its stated range.
All identities in (MA-PR.13) follow by multiplying
the full local Taylor classes, using
$\mathcal M(g*h)=\mathcal Mg\,\mathcal Mh$ and
$k_g*k_h=k_{g*h}$.  Equality of all full jets is
equality in the actual smooth quotient.
These identities are for the constructed operators;
they do not presuppose a direct-sum topology
$Q_\infty=Q_{\rm in}\oplus Q_{\rm out}$.
\end{proof}

The construction also gives the genuine covariance
$S_gU_t^{Q_H}=e^{-at}U_t^{Q_\infty}S_g$.
On every retained local algebra both sides multiply
by $e^{(\rho-a+z)t}\mathcal Mg(\rho-a+z)$, with the
full Taylor class understood.  The exterior jets
are zero.  This is another equality of actual maps,
by the full jet kernel.

\paragraph{MA-PR-003 (partial Meyer transport with an operator correction).}
The traces of the actual operators satisfy
\[
 \operatorname{Tr}_{Q_H}C_g
    =\operatorname{Tr}_{Q_\infty}B_g
    =\operatorname{Tr}_{Q_\infty}C_{k_g}
                         -\operatorname{Tr}_{Q_\infty}E_g .
\tag{MA-PR.19}
\]
Under the fixed-sector isomorphism of MT-M22/MT-M24,
the middle and last terms are traces of the corresponding
actual operators on Meyer's negative smooth quotient.
Thus the Hilbert trace has been transported to a
geometric smooth representation, with the discarded
part represented by the actual nuclear operator $E_g$.

In the rational trivial sector, specialize to $a=1/2$,
$u=-\alpha$, $v=\alpha$ satisfying (MA-PR.2).
Let $W_S(g)$ be the actual local trace in (MA-MEY.15),
and retain $B_{\overline S}(g)$ and $P(g)$ from MA-RS.
Then
\[
 \operatorname{Tr}_{Q_H}C_g
      =P(g)-W_S(g)-B_{\overline S}(g)
                            -\operatorname{Tr}_{Q_\infty}E_g .
\tag{MA-PR.20}
\]
For $g(0)=0$ the literal cutoff limit of MA-WP-003 is
\[
\begin{aligned}
 \lim_R\operatorname{Tr}(M_{\chi_R}C_gP_G)
    &=W_S(g)+B_{\overline S}(g)-P(g)
                           +\operatorname{Tr}_{Q_\infty}E_g\\
    &=A_0(g)+B(g)-P(g)
                           +\operatorname{Tr}_{Q_\infty}E_g .
\end{aligned}
\tag{MA-PR.21}
\]
The transport to (MA-PR.21) uses actual operator
factorization, trace cyclicity and Meyer's actual
common-space geometry.  After this transport has
been established, the finite-block trace calculation
identifies its correction as
\[
 \operatorname{Tr}_{Q_\infty}E_g
       =\sum_{\rho\in Z_{\rm out}}
                         m_\rho\mathcal Mg(\rho-a).
\tag{MA-PR.22}
\]
This last scalar description is not used in the
construction or in the trace cyclicity step.

\begin{proof}
Pass the target vectors of (MA-PR.17) to $Q_\infty$.
Composing that series with the continuous $i$ gives
a trace-class rank-one series on the Hilbert space.
Composing the source functionals with $i$ gives a
nuclear rank-one series for $B_g$ on $Q_\infty$.
The two rank-one traces are termwise equal:
\[
 \lambda_{n,j}\alpha_{n,j}
                       (i q_\infty u_{n,j}).
\tag{MA-PR.23}
\]
Their absolute summability follows from boundedness
of the target disk, continuity of $i$, the unit
functional bound and $\sum|\lambda_{n,j}|<\infty$.
The smooth quotient has the approximation property
directly from the absolute finite-block synthesis
of MT-M18.  Equivalently the order-zero nuclear
trace is the same common series, as in MA-HM-002.
This proves the first equality of (MA-PR.19) with
the actual identity $iS_g=C_g^{Q_H}$.
Linearity of the nuclear trace and (MA-PR.12)
prove its second equality.

The MT-M22/MT-M24 map is an actual continuous
equivariant isomorphism to the fixed compact-character
sector of Meyer's negative quotient.
The normalization check in MA-HM-002, especially
(MA-HM.15) and its proof, carries $C_{k_g}$ to the
integrated Meyer operator with test
$f(ca(x))=\beta(c)^{-1}e^{-ax}g(x)$, with compact
group Haar measure of mass one.
It also carries $B_g$ and $E_g$ to their actual
conjugate operators and preserves every rank-one
pairing in (MA-PR.23).
For the rational specialization, Meyer's source
trace additivity (MA-MEY.15)--(MA-MEY.16) expresses
the full smooth trace as $P-W_S-B_{\overline S}$.
Substitution into the operator identity (MA-PR.19),
not substitution of a zero sum, proves (MA-PR.20).

MA-WP-003 proves trace class of $C_gP_Z$ on the
specified weighted Hilbert space and, in the stated
cutoff order, identifies the cutoff limit with the
negative of the actual quotient trace when $g(0)=0$.
Applying that result to (MA-PR.20) proves the first
line of (MA-PR.21).
The checked rational local-place normalization
$W_S+B_{\overline S}=A_0+B$ gives its second line.
Only the spectral product against a cutoff is
asserted to converge in trace norm; the geometric
cutoff operators themselves need not do so.

Finally $E_g$ is zero on the retained local factors
and is multiplication by the full Taylor class of
$\mathcal Mg(w-a)$ on each exterior factor.
The actual finite-block trace is therefore
$m_\rho\mathcal Mg(\rho-a)$ on an exterior factor
and zero on an interior factor.
The absolute block nuclear series just used, or
the full and partial series of MA-WP, permits
summation of these finite-dimensional traces.
The local zero count and vertical Gevrey decay
also prove absolute convergence directly.
This proves (MA-PR.22).
\end{proof}

\paragraph{MA-PR-004 (sharp obstruction at a truncated boundary block).}
For this paragraph do not assume (MA-PR.2).
Suppose a zero $\rho$ is on an edge of the Hilbert strip and
MA-WP retains a proper nonzero part of its multiplicity:
\[
       m=m_\rho,\qquad r=r_\rho,\qquad 0<r<m,\qquad
                         \lambda=\rho-a .
\tag{MA-PR.24}
\]
The actual smooth and Hilbert primary factors at this point
are, with their translation actions,
\[
 A=\mathbb C[z]/(z^m),\qquad B=\mathbb C[z]/(z^r),
 \qquad \pi_{m,r}:A\longrightarrow B .
\tag{MA-PR.25}
\]
The restriction of $i$ is exactly this quotient map.
For a test $g$ write $G(z)=\mathcal Mg(z)$.
A translation-equivariant return of this primary factor
which factors the actual convolution must satisfy
\[
 L:B\longrightarrow A,\qquad
       \pi_{m,r}L=M_{G(\lambda+z)}^B .
\tag{MA-PR.26}
\]
Here equivariance includes the centering factor
$e^{-at}$, as for $S_g$ above.
Such a linear equivariant $L$ exists if and only if
\[
     G^{(j)}(\lambda)=0
           \quad(0\le j<\min(m-r,r)).
\tag{MA-PR.27}
\]
This is a sharp finite-block condition.  In particular,
if $G(\lambda)\ne0$, no such return exists, even before
continuity is required.

Every prescribed boundary point admits a compact
finite-order Gevrey test with $g(0)=0$ and
$G(\lambda)=1$.  Thus no translation-equivariant
return can factor convolution for all these tests
at a proper truncated boundary factor.  In particular
this factor cannot be realized as a smooth invariant
subrepresentation on which $i$ is an isomorphism.
The assertion is conditional on such a truncated
boundary factor actually occurring; it does not
assert a multiple zero of any particular $L$-function.

\begin{proof}
MA-WP-001 identifies the actual Hilbert factor with
$B$, with translation $e^{(\lambda+z)t}$.
The full smooth factor has translation
$e^{(\rho+z)t}$ by MT-M18.
After the centering factor is removed, equivariance
is
\[
                      L e^{tz}=e^{tz}L\quad(t\in\mathbb R).
\tag{MA-PR.28}
\]
These are finite polynomial matrices in $t$, so
differentiation at zero gives $Lz=zL$.
Equivalently $L$ is a $\mathbb C[z]$-module map.
It is determined by $c(z)=L(1)\in A$, and its
well-definedness on $B$ is exactly $z^r c=0$ in $A$.
Consequently
\[
\begin{aligned}
 \operatorname{Hom}_{\mathbb C[z]}(B,A)
       &\simeq \operatorname{Ann}_A(z^r)=z^{m-r}A,\\
 \pi_{m,r}\operatorname{Ann}_A(z^r)
       &=z^{m-r}B .
\end{aligned}
\tag{MA-PR.29}
\]
If $m-r\ge r$, the ideal on the second line is zero.
Otherwise it consists precisely of the classes
whose coefficients of degrees $0,\ldots,m-r-1$
vanish.  Equation (MA-PR.26), applied to $1$,
therefore has a solution exactly when
$G(\lambda+z)\bmod z^r$ belongs to that ideal.
This is (MA-PR.27).

For the converse, when $m-r<r$ take as $c(z)$ the
degree-less-than-$r$ Taylor polynomial of
$G(\lambda+z)$.  Under (MA-PR.27) it is divisible
by $z^{m-r}$, so $z^r c=0$ in $A$ and
$L(b)=bc$ gives the required module map.
When $m-r\ge r$, the condition makes the multiplier
on $B$ zero and $L=0$ suffices.
More generally all choices are
\[
              c=c_0+z^{\max(r,m-r)}A,
\tag{MA-PR.30}
\]
where $c_0$ is any one solution, because the difference
must lie in both $z^{m-r}A$ and $\ker\pi_{m,r}=z^r A$.
This proves the sharpness assertion without a
semisimplicity assumption.

To see why the finite-block condition is necessary
for a proposed global return, embed the actual
Hilbert factor $B$ in its finite block using the
primary idempotent of MA-WP, and then project the
smooth output to its full $\rho$-jet factor.
The finite-dimensional jet maps are continuous.
A globally equivariant return yields (MA-PR.28)
on these factors, and the actual identity
$iS=C_g^{Q_H}$ yields (MA-PR.26).
Thus no other block can remove this obstruction.
Conversely the finite solution just constructed
can be represented in the smooth quotient by
MT-M18's finite-block interpolation.  This converse
is local to the one primary factor; it does not
assert estimates for a return on an infinite
collection of blocks.

Choose a compact Gevrey bump $\eta$ supported in
$(1,2)$ with $\int\eta=1$, and put
$g(x)=e^{-\lambda x}\eta(x)$.
Multiplication by this analytic exponential
preserves its Gevrey order on the compact support.
Then $g(0)=0$ and $\mathcal Mg(\lambda)=1$.
The condition (MA-PR.27) fails.
Finally an invariant smooth realization of $B$
with $i$ an isomorphism would give a module section
of $\pi_{m,r}$, corresponding to the unit multiplier
in (MA-PR.26), which has just been excluded.
\end{proof}

\paragraph{MA-PR-005 (Gaussian-regularized actual returns for all smooth tests).}
Return to the generic-edge hypothesis (MA-PR.2).
For every $g\in\mathcal E$, put
\[
 \gamma_\varepsilon(x)
       =(4\pi\varepsilon)^{-1/2}e^{-x^2/(4\varepsilon)},
 \qquad g_\varepsilon=g*\gamma_\varepsilon,\qquad
                \mathcal M\gamma_\varepsilon(z)=e^{\varepsilon z^2}.
\tag{MA-PR.31}
\]
One may take $h_\varepsilon=g_\varepsilon$.
If $g(0)=0$, fix once and for all a compact finite-order
Gevrey function $b$ with $b(0)=1$, and instead take
\[
       h_\varepsilon
            =g_\varepsilon-g_\varepsilon(0)b .
\tag{MA-PR.32}
\]
Then $h_\varepsilon(0)=0$ and in either case
$h_\varepsilon\to g$ in $\mathcal E$.

For every fixed $\varepsilon>0$, (MA-PR.15) defines
an actual continuous return
$S_\varepsilon:Q_H\to Q_{\rm in}$, nuclear of every
positive order, with
\[
\begin{aligned}
 iS_\varepsilon&=C_{h_\varepsilon}^{Q_H},\\
 B_\varepsilon&=S_\varepsilon i,\qquad
 E_\varepsilon
   =C_{e^{-a\cdot}h_\varepsilon}^{Q_\infty}-B_\varepsilon .
\end{aligned}
\tag{MA-PR.33}
\]
The actual Hilbert operators satisfy
\[
       C_{h_\varepsilon}^{Q_H}
                \longrightarrow C_g^{Q_H}
                           \quad\hbox{in trace norm}.
\tag{MA-PR.34}
\]
The full smooth trace and the discarded correction traces
have the separate limits
\[
\begin{aligned}
 \operatorname{Tr}_{Q_\infty}
       C_{e^{-a\cdot}h_\varepsilon}
     &\longrightarrow
       \operatorname{Tr}_{Q_\infty}C_{e^{-a\cdot}g},\\
 \operatorname{Tr}_{Q_\infty}E_\varepsilon
     &\longrightarrow
        \operatorname{Tr}_{Q_\infty}C_{e^{-a\cdot}g}
                     -\operatorname{Tr}_{Q_H}C_g .
\end{aligned}
\tag{MA-PR.35}
\]
No convergence of $S_\varepsilon$, $B_\varepsilon$
or $E_\varepsilon$ as maps into $Q_\infty$ is
asserted for general $g$.

In particular, for every compact smooth $g$ with $g(0)=0$,
the rational centered geometric cutoff limit at this
generic width has the actual-operator transport
\[
 \lim_R\operatorname{Tr}(M_{\chi_R}C_gP_G)
     =A_0(g)+B(g)-P(g)
             +\lim_{\varepsilon\downarrow0}
                         \operatorname{Tr}_{Q_\infty}E_\varepsilon .
\tag{MA-PR.36}
\]
Its correction is the trace limit of the constructed
nuclear discarded operators.  As an additional
evaluation of that limit, for every $g\in\mathcal E$,
\[
 \lim_{\varepsilon\downarrow0}
             \operatorname{Tr}_{Q_\infty}E_\varepsilon
       =\sum_{\rho\in Z_{\rm out}}
                              m_\rho\mathcal Mg(\rho-a).
\tag{MA-PR.37}
\]

\begin{proof}
The Gaussian belongs to $\mathcal E$ for each
$\varepsilon>0$, and its Mellin transform in
(MA-PR.31) follows by the Gaussian integral and
analytic continuation in $z$.
The convolution estimates of MT-M19 give
$g_\varepsilon\in\mathcal E$ and
\[
 \mathcal Mg_\varepsilon(z)
       =\mathcal Mg(z)e^{\varepsilon z^2}.
\tag{MA-PR.38}
\]
Here the actual multiplier in the uncentered
$w$-block is
$\mathcal Mg(w-a)e^{\varepsilon(w-a)^2}$.
This fixes the centering in the regularization.

For completeness the approximation takes place in
every defining seminorm, not merely in $L^2$.
The fundamental theorem of calculus gives
\[
 p_{N,j}(U_y g-g)
       \le |y|e^{N|y|}p_{N,j+1}(g).
\tag{MA-PR.39}
\]
Integrate against $\gamma_\varepsilon(y)\,dy$.
After $y=\sqrt{\varepsilon}\,t$, its right-hand
moment is at most $C_N\sqrt{\varepsilon}$ for
$0<\varepsilon\le1$.
Thus
\[
 p_{N,j}(g_\varepsilon-g)
       \le C_N\sqrt{\varepsilon}\,p_{N,j+1}(g).
\tag{MA-PR.40}
\]
When $g(0)=0$, this also proves
$g_\varepsilon(0)\to0$, so subtraction of the
fixed $b$ in (MA-PR.32) preserves convergence
in every $\mathcal E$ seminorm.

The only use of compact Gevrey regularity in the
proof of MA-PR-002 was the multiplier estimate
(MA-PR.14).  The remainder of that proof applies
to any $h\in\mathcal E$ whose contour multipliers
are bounded by $C e^{-c\nu_n^\theta}$ for some
$\theta>0$: (MA-PR.16) is summable in every
seminorm, the square-root construction (MA-PR.17)
gives the common nuclear series, and the actual
jet identities give (MA-PR.18).
For $g_\varepsilon$, (MA-PR.38) on the fixed
real strip gives precisely such a bound:
\[
 \sup_{\partial\Omega_n^{\rm in}}
       |\mathcal Mg_\varepsilon(w-a)|
            \le C_\varepsilon e^{-c_\varepsilon\nu_n^2}.
\tag{MA-PR.41}
\]
The unregularized transform $\mathcal Mg$ is
bounded on that fixed strip, and the real part
of $\varepsilon(w-a)^2$ is
$\varepsilon((\operatorname{Re}w-a)^2-(\operatorname{Im}w)^2)$.
This proves (MA-PR.41) with adjusted constants
for the finitely many low blocks.
If (MA-PR.32) is used, the additional compact
Gevrey transform has a bound
$C e^{-c\nu_n^{1/q}}$.
The sum still has the required positive-power
exponential decay for each fixed $\varepsilon$.
No uniform estimate as $\varepsilon\downarrow0$
is needed or claimed for these return maps.
The same construction therefore proves
(MA-PR.33), together with its actual full-jet
descriptions and nuclear trace cyclicity.

We now pass limits in two different spaces.
MA-WP.14--MA-WP.15 gives a continuous map
$\mathcal E\to\mathcal S_1(Q_H)$.
More explicitly, summing its rank-times-operator-norm
bound after a sufficiently strong vertical-decay
estimate bounds
$\|C_f^{Q_H}\|_1$ by finitely many
$\mathcal E$ seminorms of $f$.
Apply this to $f=h_\varepsilon-g$ and (MA-PR.40).
This proves the Hilbert trace-norm convergence
(MA-PR.34).

For the full smooth term use the fixed block
series of MT-M19, (M19.3)--(M19.4), on
$Q_\infty\simeq s(J)$.
Its Auerbach coefficient functionals are
equicontinuous, its coefficients
$(2+|n|)^{-2}$ are summable with the block
dimensions, and the target vectors depend
continuously on the test in every bounded-family
seminorm.
For the difference $f=e^{-a\cdot}(h_\varepsilon-g)$,
the sum of absolute rank-one trace pairings is
at most
\[
       \sum_n d_n\,
          \sup_{w\in\partial\Omega_n}|\mathcal Mf(w)|
                \longrightarrow0 .
\tag{MA-PR.42}
\]
Indeed integration by parts on the fixed full
strip bounds the supremum by $\nu_n^{-2}$
times finitely many seminorms of $f$, and
$d_n=O(\log\nu_n)$.
This proves the first limit in (MA-PR.35)
from the actual smooth nuclear series.

For each fixed $\varepsilon$, the common series
of the already constructed return gives
\[
 \operatorname{Tr}_{Q_H}C_{h_\varepsilon}
    =\operatorname{Tr}_{Q_\infty}
            C_{e^{-a\cdot}h_\varepsilon}
                   -\operatorname{Tr}_{Q_\infty}E_\varepsilon .
\tag{MA-PR.43}
\]
The first trace converges by (MA-PR.34), and
the full smooth trace converges by (MA-PR.42).
Their difference proves the second limit of
(MA-PR.35).  It is a statement about correction
traces, without an operator convergence claim
for $E_\varepsilon$ on the smooth quotient.

For the rational centered compact test $g$,
apply Meyer's common-space formula to the
unregularized full smooth operator
$C_{e^{-x/2}g}$, whose trace is the limit of the
full smooth terms just proved.
It gives $P-W_S-B_{\overline S}$ as before.
Apply the actual MA-WP-003 ambient cutoff
theorem to $g$ itself, giving the negative
of the Hilbert trace.
Equations (MA-PR.35) and these two actual
identifications prove (MA-PR.36).
This argument does not require a new
common-space geometric formula or a compact
support assertion for the Gaussian-regularized
tests.

Lastly the constructed $E_\varepsilon$ has
zero retained jets and the full exterior
Taylor multiplier of $h_\varepsilon$.
The finite-block calculation of MA-PR-003
therefore evaluates its trace as the
exterior divisor sum for $h_\varepsilon$.
Convergence $h_\varepsilon\to g$ in
$\mathcal E$, the local zero count and
rapid vertical Mellin decay give convergence
of these sums absolutely.
This proves (MA-PR.37) after the actual
operator and trace limits have been established.
\end{proof}

\paragraph{Scope, boundary and the remaining independent-method step.}
The maps (MA-PR.11)--(MA-PR.12) are unconditional under
the explicit generic-edge condition and Gevrey test
hypothesis.  They apply at every fixed polynomial
weight $s\ge0$, including the requested $s\ge3$.
No boundary jets occur in those unregularized maps
because (MA-PR.2) excludes boundary zeros.
The exact boundary truncations at other widths
remain those of MA-WP.  MA-PR-004 proves a sharp
local obstruction to an equivariant return when
a nonzero proper part of a boundary block survives;
it gives the local converse for tests satisfying
the necessary vanishing condition, without
asserting a global return at exceptional widths.
MA-PR-005 gives a family of actual regularized
returns for arbitrary $\mathcal E$ tests at
generic widths, with Hilbert trace-norm convergence
and correction-trace convergence.  It does not
remove the Gevrey hypothesis for an unregularized
map $S_g$ into $Q_\infty$.

If exterior zeros occur, $S_gi$ is the retained
operator $B_g$, not the full convolution $C_{k_g}$.
MA-HM-003 excludes the latter identity for a test
that detects a discarded block.
Nothing above proves an unsmoothed continuous
spectral selector, nor a full return on the same
partial completion.  The identity (MA-PR.21)
therefore retains its actual discarded-divisor
operator.  Removing that term for every test is
the exact criterion of MA-EG-001.

The paired regularized Euler diagonals of MA-RS
carry the same discarded-divisor trace in their
contour correction.  That agreement does not
identify the unbounded inverse composition there
with every off-diagonal map in Meyer's original
common-space projection.  The factorization
proved here is the explicit map (MA-PR.15) through
the actual retained smooth subrepresentation;
these are distinct domain statements.

The Grothendieck placement concerns actual finite-length
$\mathbb C[z]$-modules at the retained and discarded
spectral points.  Their local lengths and underlying
finite-dimensional complex endomorphisms enter the
nuclear traces in (MA-PR.19)--(MA-PR.23).
This is the same inherited cote~1, scans~079--080
and~111--112, placement as MA-HM and MA-WP.
The analytic factorization does not construct an
arithmetic positivity or weight realization.

\paragraph{Exact inputs and proof scope.}
The arithmetic estimates are the unconditional local
zero count and fixed-strip finite-$L$ bounds of NT-Z10,
and the entire lift, selected heights, Stirling bounds
and full smooth zero-jet kernel of MT-M18,
(M18.5)--(M18.13).  The new inner-contour lower bound
and all restriction and nuclear estimates are proved
above.  MA-WP-001 supplies the already proved literal
Hilbert closure and its finite retained blocks;
MA-WP-003 supplies the actual ambient cutoff limit.
The Meyer identifications and common-space signs
are exactly MT-M22/MT-M24, MA-HM-002 and
MA-MEY.15--MA-MEY.16.  No A, M, or new outside
preprint is an input to this generic partial return.
The boundary obstruction uses only the exact
local modules of MA-WP and elementary finite
$\mathbb C[z]$-module algebra.  The Gaussian
extension uses the explicit approximate-identity
estimate (MA-PR.40), the literal Hilbert
trace-norm continuity in MA-WP, and the full
smooth rank-one continuity in MT-M19.
